\subsection{Alternative Designs on Energy Functions}
\label{section:alternative designs}
Our proposed energy functions $E_\text{ATTN}$ and $E_\text{FF}$ provide an avenue to quantify the objective on uniformity and alignment in an amenable way for optimization. To manifest the significance of our conceptualization in designing a variety of Transformer-based models, we extend the energy functions to more general forms and provide alternative instantiations of them that can induce novel structures. Specifically, we generalize the energy functions in~\Eqref{eq:2} and~\Eqref{eq:4} to the following forms:
\begin{align}
E_\text{ATTN} &= \sum_{h=1}^H E_\text{ATTN}^h = \sum_{h=1}^H\sum_{i=1}^N f\left( \sum_{j=1}^N K(\mW_h^\top\vx_i, \mW_h^\top\vx_j)\right), \label{eq:generalized attn} \\
E_\text{FF} &= - \sum_{i=1}^N g\left(\sum_{m=1}^M h(\vd_m^\top \vx_i)\right). \label{eq:generalized ff}
\end{align}
where $K:\mathbb{R}^p \times \mathbb{R}^p \rightarrow \mathbb{R} $ is a kernel function and $f,g,h: \mathbb{R} \rightarrow\mathbb{R}$ are non-decreasing scalar functions. For clarity, we omit the hyperspherical constraint but its correspondence to $\operatorname{RMSNorm}$ still holds. 

By choosing different variations on these functions, we arrive at alternative energy functions and their consequent attention and feedforward operators, of which we summarize the specifications in the following tables. Remarkably, in Table~\ref{tab:generalized attention}, we can derive an attention with linear complexity $\mathcal{O}(N)$ by specifying the kernel function with the inner product of an element-wise transformation $\Phi: \mathbb{R} \rightarrow \mathbb{R}$, bridging our energy view with recent advances in linear attention design. In practice, we choose it as the sigmoid function $\sigma(x) = 1/(1+\exp(-x))$, but other designs can also be possible. In Table~\ref{tab:generalized feedforward}, if we specify the outer function $g$ in the feedforward energy as a quadratic function, there is a novel summation and a Hadamard product operation emerging with the transformation $\Phi$, similar to the gating mechanism. We also specify $\Phi$ with the sigmoid function.


In summary, these design choices demonstrate that \textsc{Hyper-SET} is more than just a single model, but a blueprint for constructing principled Transformer variants. Each component—self-attention, feedforward, and normalization—can be systematically interpreted and designed within our energy minimization framework, providing a pathway for principled, modular innovation in sequence model architectures. 

\begin{table}[htbp]
\caption{Alternative designs on attention energy $E_\text{ATTN}$ and the induced operators.}
\label{tab:generalized attention}
\centering
\resizebox{\linewidth}{!}{
\begin{tabular}{lllll}
% \begin{tabularx}{\textwidth}{l*{4}{X}}
\toprule
% \multirow{2}{*}{Models}  & \multirow{2}{*}{Config ($\#$ Params)} & \multicolumn{3}{c}{Dataset}\\
% \cmidrule(lr){3-5}
%  &  & CIFAR-10 & CIFAR-100 & IM-100\\
\textbf{Operator }
& $f(x)$ 
& $K(\bm{x},\bm{y})$ 
& $E_\text{ATTN}$ 
& $-\nabla_{\mX} E_\text{ATTN}$ \\
\midrule
Bi-Softmax (Default) 
& $\beta^{-1}\log(x)$ 
& $\exp(\beta \bm{x}^\top\bm{y})$ 
& \Eqref{eq:2} 
& \Eqref{eq:6} \\
Sigmoid Attention  
& $\frac{\beta^{-1}}{2}x$ 
& $\sigma(\beta \bm{x}^\top\bm{y})$ 
& $\frac{1}{2}\sum\limits_{h=1}^H\sum\limits_{i,j=1}^N\sigma\left(\beta (\mW_h^\top\vx_i)^\top\mW_h^\top\vx_j\right) \beta^{-1}$ & $\sum\limits_{h=1}^H \mW_h \mW_h^\top \mX \sigma(1-\sigma)\left( \beta (\mW_h^\top \mX)^\top \mW_h^\top \mX\right)$ \\
Linear Attention 
& $\frac{\beta^{-1}}{2}x$ 
& $\frac{1}{2}\left(\beta\Phi(\bm{x})^\top\Phi(\bm{y})\right)^2$ 
& $\frac{1}{4}\sum\limits_{h=1}^H\sum\limits_{i,j=1}^N\left( \beta \Phi(\mW_h^\top\vx_i)^\top \Phi(\mW_h^\top\vx_j)\right)^2 \beta^{-1}$ 
& $\sum\limits_{h=1}^H \mW_h \Phi^{\prime}(\mW_h^\top\mX) \odot \left(\beta \Phi(\mW_h^\top\mX)\Phi(\mW_h^\top\mX)^\top\Phi(\mW_h^\top\mX)\right)$ \\
\bottomrule
% \end{tabularx}
\end{tabular}
}
\end{table}

\begin{table}[htbp]
\caption{Alternative designs on feedforward energy $E_\text{FF}$ and the induced operators.}
\label{tab:generalized feedforward}
\centering
\resizebox{\linewidth}{!}{
\begin{tabular}{lllll}
% \begin{tabularx}{\textwidth}{lXlll}
\toprule
% \multirow{2}{*}{Models}  & \multirow{2}{*}{Config ($\#$ Params)} & \multicolumn{3}{c}{Dataset}\\
% \cmidrule(lr){3-5}
%  &  & CIFAR-10 & CIFAR-100 & IM-100\\
\textbf{Operator}
& $g(x)$ 
& $h(x)$ 
& $E_\text{FF}$ 
& $-\nabla_{\mX} E_\text{FF}$ \\
\midrule
ReLU FF (Default) 
& $x$ 
& $\frac{1}{2}\operatorname{ReLU}^2(x)$ 
& \Eqref{eq:4} 
& \Eqref{eq:9} \\
Softmax FF 
& $\log(x)$ 
& $\exp(x)$ 
& $-\sum\limits_{i=1}^N\log\left(\sum\limits_{m=1}^M\exp(\vd_m^\top\vx_i) \right)$ 
& $\mD \underbrace{\operatorname{softmax}}_{\text{column-wise}}(\mD^\top \mX)$ \\
Gated FF 
& $\frac{1}{2}x^2$ 
& $\Phi(x)$ 
& $-\frac{1}{2} \sum\limits_{i=1}^N\left(\sum\limits_{m=1}^M \Phi\left(\vd_m^\top \vx_i \right) \right)^2$ 
&  $\mD \underbrace{\Phi(\mD^\top\mX)}_{\text{column sum}} \odot ~\Phi^{\prime}(\mD^\top\mX)$\\
\bottomrule
% \end{tabularx}
\end{tabular}
}
\end{table}

\subsection{Preliminary Results on Scalability}
\label{section:scalability}
% \todo{This design maintains the core recurrence-based parameter sharing of \textsc{Hyper-SET} while allowing for depth-aware expressiveness, enabling better adaptation to increasingly complex tasks as the model scales. ; both AB are initialized as Gaussain with 0.02 std; scaling factor 4}
\paragraph{Setups for Depth-wise Flexible Computation.} To equip our model with flexible computation while maintaining its core recurrence-based parameter sharing, we add an independent low-rank adaptation \citep{hu2022lora} matrix $W=AB$ to every iteration of the Transformer layer sharing the same base parameters. Inspired by but unlike depth-wise adaptation of pre-trained models in \citet{bae2025relaxed}, we train the base parameters and the adaptation matrix together. Both matrix $A \in \mathbb{R}^{d\times r},B \in \mathbb{R}^{r\times d}$ with rank $r$ are initialized with Gaussain of 0.02 standard deviation. The scaling factor before the adaptation matrix is set as 4.

\paragraph{Scaling Results on Images.} To further demonstrate the capability of our model when scaling up its size, we perform a preliminary evaluation on image classification. Following the experimental setups in the main paper to configure the models with one layer and repeat with 12 iterations, we scale up the width 
 of Transformer to 1152 and ours to 1536, resulting in similar parameter size. The result is presented in Table~\ref{tab:scaling up}. These explorations showcase the potential of \textsc{Hyper-SET} in improving model capacity without significantly increasing the total parameter count.

\begin{table}[htbp]
\caption{Top-1 accuracy for image classification when scaling up the model size. Our model surpasses other baselines while being parameter-efficient compared to vanilla Transformer.}
\label{tab:scaling up}
\centering
\resizebox{\linewidth}{!}{
\begin{tabular}{lllcc}
\toprule
% \multirow{2}{*}{Models}  & \multirow{2}{*}{Config ($\#$ Params)} & \multicolumn{3}{c}{Dataset}\\
% \cmidrule(lr){3-5}
%  &  & CIFAR-10 & CIFAR-100 & IM-100\\
\textbf{Dataset} & \textbf{Model} & \textbf{Width} $d$ & \textbf{$\#$ Params (M)} & \textbf{Accuracy ($\%$)} \\
\midrule
\multirow{4}{*}{CIFAR-10} & Transformer & 1152 & 16.0  & 89.42 \\
                          & \textsc{CRATE} \citep{yu2023white} & 1152 & 4.1  & 88.77 \\
                          & Energy Transformer \citep{hoover2024energy} & 1152 & 8.0  & 76.21 \\
                          & \textsc{Hyper-SET} (Ours) & 1536 & 14.3  & \textbf{90.62} \\
\midrule
\multirow{4}{*}{CIFAR-100} & Transformer & 1152 & 16.1  & 62.83 \\
                          & \textsc{CRATE} \citep{yu2023white} & 1152 & 4.2  & 63.39 \\
                          & Energy Transformer \citep{hoover2024energy} & 1152 & 8.1  & 55.47 \\
                          & \textsc{Hyper-SET} (Ours) & 1536 & 14.4  & \textbf{66.30} \\
\bottomrule
\end{tabular}
}
\end{table}


\paragraph{Scaling Results on Texts.} We also provide the results for text classifications to show the potential of \textsc{Hyper-SET}. We use \texttt{bert-base-uncased} from Hugging Face as the tokenizer with a maximum sequence length of 128. The number of recurrence is set to 6 (L1R6) and width 
 is set to 384. The training recipe is the same as the image classification in Table~\ref{tab:training recipe for image classfication} except we train for 10 epochs. The results on \texttt{yelp\_review\_full} dataset, a 5-way classification task, are shown in Table~\ref{tab:text clasification}.
 \begin{table}[htbp]
\caption{Top-1 accuracy for text classification on \texttt{yelp\_review\_full} dataset.}
\label{tab:text clasification}
\centering
\begin{tabular}{llcc}
\toprule
\textbf{Model} & \textbf{Width} $d$ & \textbf{$\#$ Params (M)}& \textbf{Accuracy ($\%$)} \\
\midrule
Transformer & 384 & 13.49 	& 60.09 \\
\textsc{CRATE}-T \citep{hu2024an} & 384 & 12.02 & 59.43 \\
\textsc{CRATE} \citep{yu2023white} & 384 & 12.17  & 59.18 \\
Energy Transformer \citep{hoover2024energy} & 384 & 12.61  & 54.93 \\
\textsc{Hyper-SET} (Ours) & 384 & 12.66 & \textbf{60.75} \\
\bottomrule
\end{tabular}
\end{table}